\documentclass[11pt]{article}
\usepackage{fontspec}
\usepackage[french,english]{babel}
\usepackage{amsmath}
\usepackage{amssymb}
\usepackage{amsthm}
\usepackage{mathrsfs}
\usepackage{tikz-cd}
\usepackage{enumitem}
\usetikzlibrary{babel}

\newcommand{\printedpage}[1]{\marginpar{\footnotesize p.~#1}}
\newcommand{\msp}[1]{\marginpar{\footnotesize [#1]}}
\newcommand{\circled}[1]{\tikz[baseline=(c.base)]\node[draw,circle,inner sep=1pt](c){#1};}
\newcommand{\Mc}{\mathcal{M}}
\newcommand{\U}{\mathfrak{U}}

\begin{document}

\printedpage{49}
% head: A. Grothendieck : Dérivateurs, Ch. XVIII, Édition des Universités de Montpellier II et Paris VII
% foot: Version 17 décembre 2025
% The mark of footnote 59 is on the preceding page of the edition; its text is printed at the foot of this page.
\footnotetext[59]{\textbf{Question.} Condition $e_{\pi_0}$ (or even, the pair $\alpha(\pi_0) + \beta(\pi_0)$ plus the existence of a small strictly generating subcategory (and the stability of $\Mc$ under small $\varinjlim$ ), does it imply that $\Mc$ is a $\pi_0$-paratopos ? I presume not.)}

\begin{itemize}[leftmargin=4em]
\item[\circled{e}] (\footnote[60]{This is the prettiest intrinsic characterization of paratopoi in the ``economy'' mode, the other, the structural statement, is \circled{c} .}). \textit{$\Mc$ is stable under small $\varinjlim$ , contains a small strictly generating full subcategory, and there exists an infinite cardinal $\pi_0$ such that condition $\mathrm{e}_{\pi_0}$ of 2°) is satisfied.}
\item[\circled{e$'$}] \textit{As e, but requiring, instead of $\mathrm{e}_{\pi_0}$, the weaker condition conjunction of $\alpha(\pi_0)$ and $\beta(\pi_0)$ in 2°.}
\end{itemize}

\msp{page 151}

\textsc{Proof}. Part 1° has already been proved. The first assertions in 2°), namely the stability of $\Mc$ under small $\varprojlim$ , is well known (follows from stability under small $\varinjlim$ , and from the existence of a small generating family, cf. SGA 4 I 8.12.7, dual version of the marginal annotation). The validity of condition $\mathrm{e}_{\pi_0}$ follows from the fact that it holds in $C^{\wedge}$ (where $C$ is the category mentioned in 1°, $\mathrm{c}_{\pi_0}$), and from the fact that the $\varprojlim$ in $\Mc$, and the $\varinjlim$ large relative to $\pi_0$ in $\Mc$, are computed in $C^{\wedge}$. We moreover find in the same way : % [sic] assertions -> assertion, singular subject "La première"

\textbf{Corollary 4.6.4.} \textit{If $\Mc$ is a $\U$-$\pi_0$-paratopos category, then for every cardinal $\pi \geqslant \pi_0$, and every category $J$ such that} $\mathrm{card}\, J \leqslant \pi$\textit{, the functor} % [sic] tout catégorie -> toute catégorie, gender agreement
\[
\underline{\mathrm{Hom}}(J,\Mc) \xrightarrow{\ \varprojlim\limits_{J}\ } \Mc
\]
\textit{is $\pi$-accessible, or what amounts to the same thing, for every small category $I$ large relative to $\pi$, the functor}
\[
\underline{\mathrm{Hom}}(I,M) \xrightarrow{\ \varinjlim\limits_{I}\ } \Mc % [sic] M -> \mathcal{M}, the first display has \mathcal{M}
\]
\textit{commutes with the $\varprojlim$ of cardinality $\leqslant \pi$.}

\msp{page 152}

On the other hand, the fact that $\mathrm{e}_{\pi_0}$ implies $\alpha(\pi_0)$ and $\beta(\pi_0)$ is immediate.

Part 3°), as regards the equivalence of conditions a) to d), is contained in 1°, and the equivalence with d) has also been seen (the only non-tautological thing being that condition c), namely $\Mc \simeq \mathrm{Ind}_{\pi_0}(C)$ with $C$ small and stable under $\varinjlim$ of cardinality $\leqslant \pi_0$, implies the existence of a cardinal filtration, which was seen in 4.5.9, and reaffirmed in 4.6.2 2° above). On the other hand, that these conditions a) to d) imply \circled{e} , and that \circled{e} implies e$'$, was seen in 2°. So it remains to see that the conditions of e$'$ imply a), and for this, it remains to establish that the elements of $\Mc$ are accessible. But this is what is established in SGA 4 I 9.11, except that the statement of condition c), second case, is botched, and should be replaced by condition $\beta(\pi_0)$ (condition $\alpha(\pi_0)$ appearing in condition a) of loc. cit.).

\msp{page 153}

\textbf{Remark 4.6.5.} The developments of sections 4.4, 4.5, 4.6 considerably sharpen and strengthen, in many respects, the theory of cardinal filtrations of SGA 4 I 9. This is so, at least, in the case of $\U$-categories $\Mc$ stable under small $\varinjlim$ . In loc. cit. one works with slightly more general $\U$-categories, those satisfying
\printedpage{50}
% head: A. Grothendieck : Dérivateurs, Ch. XVIII, Édition des Universités de Montpellier II et Paris VII
% foot: Version 17 décembre 2025
the \textit{condition} $L$ : there exists an infinite cardinal $\pi_0$ such that $\Mc$ is stable under small $\varinjlim$ \textit{large relative to} $\pi_0$. I presume however that one should be able to improve in a similar way the theory of \textit{accessible} cardinal filtrations, taking inspiration from corollary 9.18 of loc. cit., establishing an equivalence of categories
\[
\Mc \;\simeq\; \mathrm{Ind}_{\pi_0}(C)\;,
\]
for a cardinal $\pi_0$ and a suitable small category $C$. Here, one can no longer assume $C$ stable under this or that type of $\varinjlim$ . Nevertheless, 4.4.15 tells us that without any further hypothesis on $C$, $\Mc \overset{\text{def}}{=} \mathrm{Ind}_{\pi_0}(C)$ does satisfy $\mathrm{L}_{\pi_0}$, i.e. is stable (and even stable in $C^{\wedge}$)

\msp{page 154}

under small $\varinjlim$ large relative to $\pi_0$.

This seen, I feel like calling \textit{accessible} a category $\Mc$ which admits a cardinal filtration, and which is moreover weakly accessible \texttt{[i.e.]} every object of which is accessible. And to venture a % [sic] filtration cardinal -> filtration cardinale, gender agreement

\textbf{Conjecture 4.6.6} (\footnote[61]{All this will be made considerably more precise in the theory which will be developed in sections 4.7, 4.16.}). \textit{Let $\Mc$ be a $\U$-category. Equivalent conditions :} % [sic] (in footnote 61) développé -> développée, agreement with "la théorie"
\begin{itemize}[leftmargin=4em]
\item[\textit{a)}] \textit{$\Mc$ satisfies condition $L$, and admits a small strictly generating category consisting of accessible objects.}
% printed: a vertical double arrow $\Uparrow$ at the left, between a) and b)
\item[\llap{$\Uparrow$\ }\textit{b)}] \textit{There exists a small category $C$, and an infinite cardinal $\pi_0$, such that $\Mc$ is equivalent to the category $\mathrm{Ind}'_{\pi_0}(C)$.}
% printed: a vertical double arrow $\Uparrow$ at the left, between b) and c)
\item[\llap{$\Uparrow$\ }\textit{c)}] \textit{$\Mc$ is ``accessible'', i.e. is weakly accessible (every object accessible) and admits a cardinal filtration.}
\end{itemize}
We have b $\Longrightarrow$ a $\Longleftarrow$ c tautologically, and c $\Longrightarrow$ b was seen in loc. cit. 9.18, so we have the sequence
\[
\mathrm{c} \;\Longrightarrow\; \mathrm{b} \;\Longrightarrow\; \mathrm{a}\;,
\]
so it would remain to establish a $\Longrightarrow$ c.

\msp{page 155}

The work accomplished (and the sweat shed) in 4.5 will moreover save us, when the time comes, from quite natural errors in the description of a cardinal filtration, starting from a strictly generating subcategory $S$ : it is neither $\mathrm{Filt}^{\pi}$, nor $\mathrm{Filt}'^{\pi}$, but something subtler -- perhaps one must take the category $\Mc_{\pi}$ again, or else saturate $S$ under the direct limits of cardinality $\leqslant \pi$ which are willing to be representable in $\Mc$ : we then find a $\mathrm{Filt}''^{\pi} \subset \Mc_{\pi}$, and if these $\mathrm{Filt}''^{\pi}$ do indeed give a cardinal filtration, we must necessarily have $\Mc_{\pi} \subset \mathrm{Filt}''^{\pi}$, hence finally $\mathrm{Filt}''^{\pi} = \Mc_{\pi}$. So I expect that in any event, the cardinal filtration (if it exists) will be given by the $\Mc_{\pi}$. Condition a) of loc. cit. 9.12 ($\Mc_{\pi}$ equivalent to a small category), and the bound on cardinals in loc. cit. c), should not cause any difficulties, and condition b) is OK. The difficulty

\msp{page 156}

comes from the fact that $\Mc_{\pi}$ not even being stable under sums and cokernels of double arrows (when these are not represented in $\Mc$), the category $\Mc_{\pi}/x$ (for $x$ in $\Mc$) has no reason to be filtering, as far as I can see, nor a fortiori
\printedpage{51}
% head: A. Grothendieck : Dérivateurs, Ch. XVIII, Édition des Universités de Montpellier II et Paris VII
% foot: Version 17 décembre 2025
to be large relative to $\pi$. So it is not too clear how then to represent $x$ as $\varinjlim\limits_{I} x_i$, the $x_i$ in $\Mc_{\pi}$ and $I$ large relative to $\pi$.

It must indeed be said that in SGA 4 I 9, I give no \textit{example} of a cardinal filtration on a category where arbitrary small $\varinjlim$ would not be representable (cf. loc. cit. 9.13). Perhaps after all it is against nature to want them beyond ? ? Perhaps one must wait until forced to it, to try to develop a theory outside the framework of paratopoi.

\msp{page 157}

\textbf{4.6.7.} We are now in a position to reformulate the litanies of somewhat technical conditions a) to f) of theorem 2.6 (p. 25 in XVIII) concerning the ``special on the left'' character of the set $\mathrm{Mon}(\Mc) \subset \mathrm{Fl}(\Mc)$ :

\textbf{Proposition 4.6.7} (\footnote[62]{Perhaps one could call these categories ``\textit{parabelian paratopos}'' categories ? ?}). \textit{The preliminary conditions on the $\U$-category $\Mc$, ensuring the validity of theorem XVIII 2.6, in a very slightly modified form (namely by modifying in 2.6 d) the condition of stability of} $\mathrm{Filt}^{\pi}$ \textit{under quotients, into an exactness property d) below), is equivalent to the set of the following six conditions a) to f).} % [sic] modifié -> modifiée, agreement with "forme"; équivaut -> équivalent, plural subject "Les conditions"
\[
\left\{
\begin{array}{ll}
\textit{a)} & \Mc\ \textit{is stable under small}\ \varinjlim\ . \\[1.5ex]
\textit{b)} & \Mc\ \textit{admits a small strictly generating subcategory.} \\[1.5ex]
\textit{c)} & \textit{Filtering direct limits in}\ \Mc\ \textit{are exact.}
\end{array}
\right.
\]
\textit{These conditions imply that $\Mc$ is a} $\U$-paratopos category \textit{(cf. 4.6.3, 3°), and in particular that the objects of $\Mc$ are} accessible, \textit{and that $\Mc$ admits a cardinal filtration $(\Mc_{\pi})_{\pi \geqslant \pi_0}$ such that for every $\pi \geqslant \pi_0$, $\Mc_{\pi}$ is stable under $\varinjlim$ of cardinality $\leqslant \pi$ (in fact, we will have $\Mc_{\pi} =$ full subcategory of the $\pi$-accessible objects of $\Mc$). Also, they imply the stability of $\Mc$ under small $\varprojlim$ . On the other hand, the following three conditions are stability conditions for monomorphisms.}
\begin{itemize}[leftmargin=4em]
\item[\textit{d)}] \textit{Let $f : X \longrightarrow Y$ be an arrow in $\Mc$ and $R = X \times_Y X \rightrightarrows X$, then $X/R \longrightarrow Y$ is a monomorphism -- in other words, $f$ factors as a composite $ip$, with $i$ monomorphic and $p$ an effective epimorphism.}
\end{itemize}

\msp{page 158}

\begin{itemize}[leftmargin=4em]
\item[\textit{e)}] \textit{If $A_1$, $A_2$ are two subobjects of an object $B$ of $\Mc$, and $A_0 = A_1 \cap A_2 \overset{\text{def}}{=} A_1 \times_B A_2$, then $A_1 \amalg_{A_0} A_2 \longrightarrow B$ is a monomorphism. (Note the close kinship of conditions d) and e).)}
\item[\textit{f)}] \textit{The set} $\mathrm{Mon}(\Mc)$ \textit{of monomorphisms in $\Mc$ is stable under cobase change.}
\end{itemize}
In fact, it is not clear that these conditions, and in particular d), suffice to entail the stability of the $\Mc_{\pi}$ under quotients (only under effective quotients), as postulated in 2.6 d). But in the statement and the proof given in XV cor. 7.4 (p. 98-99), we were content with a somewhat weaker form of this stability property, restricting ourselves to the quotients $Z'$ of $Z \in \mathrm{Ob}\,\Mc_{\pi}$ \textit{such} that there is no subobject $Z'_0$ of $Z'$, distinct from $Z'$, and such that $Z \longrightarrow Z'$ factors through $Z'_0$ ; we therefore required that $Z' \in \mathrm{Ob}\,\Mc_{\pi}$. But if condition d) is satisfied, it follows at once that a quotient $Z'$ of $Z$ having this additional property is necessarily effective, from which it does follow that $Z' \in \mathrm{Ob}\,\Mc_{\pi}$.

\printedpage{52}
% head: A. Grothendieck : Dérivateurs, Ch. XVIII, Édition des Universités de Montpellier II et Paris VII
% foot: Version 17 décembre 2025
\msp{page 159}

\textbf{Corollary 4.6.8.} \textit{Conditions d), e), f) of 4.6.7 are satisfied if $\Mc$ is an abelian category. So the set of conditions a) to f) is satisfied in each of the following two cases :}
\begin{itemize}[leftmargin=4em]
\item[\textit{1)}] \textit{$\Mc$ is an abelian category where filtering direct limits are representable and exact, and admitting a small strictly generating category} (\footnote[63]{This is the case studied in ``Tohoku''.}).
\item[\textit{2)}] \textit{$\Mc$ is a topos.}
\end{itemize}
\textbf{Remark 4.6.9.} I am not sure I know other examples where the conditions of 4.6.7 are verified, than those noted in 4.6.8 above (\footnote[64]{or which follow trivially from them (products etc.)}). One will note that case 1°) mentioned there is that of Tohoku, where I establish the existence of ``enough injective objects'' in such an abelian category. In Tohoku, I confine myself to assuming the existence of a small subcategory, \textit{generating for monomorphisms}, a condition apparently weaker than the one stated in 4.6.8 1°). But in an abelian category, the two notions (and their various variants in SGA 4 I 7.1) are equivalent by virtue of SGA 4 I 7.3 (i). % [sic] par sûr -> pas sûr, negation "ne ... pas"

\msp{page 160}

On the other hand, it seems that conditions e) f) are practically never verified for the categories defined by the species of algebraic structure ``definable by finite products'' and nothing more, such as groups, rings (of all kinds), etc. But the category of \textit{commutative} groups, being abelian, is an exception. (One can moreover interpret it as $\Mc \boxtimes \Mc$, where $\Mc = (\mathrm{Gr})$, and one can then wonder whether there are some generalizations to paratopos categories of this form, $\Mc$ being a paratopos category \ldots) % [sic] tel les groupes -> tels les groupes, agreement; quelque généralisations -> quelques généralisations, plural

\msp{page 161}

\section*{4.7\quad $\mathrm{Ind}_{\pi}C$ ($C$ arbitrary) and generalized cardinal filtrations.}

I feel the need to launch into developments on this subject, to arrive at a better understanding of the notion of the \textit{accessible part} of a paratopos, and of the stability theorem for paratopoi under certain key operations on them.

In what follows, $C$ denotes a category equivalent to a small category, $C_0$ a skeletal exhaustive subcategory, $\pi_0$ an infinite cardinal. We set
\[
\Mc \;=\; \mathrm{Ind}_{\pi_0}C \;\subset\; C^{\wedge} \qquad \text{(\footnotemark[65])}\;. \tag*{(4.7.1)}
\]
\footnotetext[65]{\textbf{N.B.} I now denote by $\mathrm{Ind}C$, $\mathrm{Ind}_{\pi}C$ what was previously denoted $\mathrm{Ind}'C$, $\mathrm{Ind}'_{\pi}C$.}
We propose to introduce on $\Mc$ various $\pm$ ``cardinal'' filtrations, by corresponding properties of an object $F$ of $\Mc \subset C^{\wedge}$, which I am now going to introduce, taking inspiration from 4.5.2 (p. 107) and 4.5.6 (p. 116).

\textbf{Proposition 4.7.1.} \textit{Set} $\pi_1 = \sup(\pi_0, \mathrm{card}\, C_0)$. \textit{Let $\pi \geqslant \pi_0$ be a cardinal, and let us consider the following conditions (depending on $\pi$) on an object $F$ of $\Mc$ :}
\begin{itemize}[leftmargin=4em]
\item[\textit{a)}] $\mathrm{card}\, C_0/F \leqslant \pi$ \textit{(i.e.} $\mathrm{card}\,\mathrm{Fl}\, C_0/F \leqslant \pi$\textit{).}
\item[\textit{a$'$)}] $\mathrm{card}\,\mathrm{Ob}\, C_0/F \leqslant \pi$.
\item[\textit{a$''$)}] $\mathrm{card}\, F(x) \leqslant \pi$ \textit{for every $x \in$} $\mathrm{Ob}\, C$ \textit{(}\textbf{N.B.} \textit{it suffices to require this for $x$ in $C_0$).}
\end{itemize}

\printedpage{53}
% head: A. Grothendieck : Dérivateurs, Ch. XVIII, Édition des Universités de Montpellier II et Paris VII
% foot: Version 17 décembre 2025
\msp{page 162}

\begin{itemize}[leftmargin=4em]
\item[\textit{b)}] \textit{We have $F = \varinjlim\limits_{i \in I} x_i$, with the $x_i \in$} $\mathrm{Ob}\, C$\textit{, $I$ a category (not necessarily ordered) large relative to $\pi_0$, and such that} $\mathrm{card}\, I \leqslant \pi$. \textit{(}\textbf{N.B.} \textit{by virtue of 4.4.2 (p.79), for given $I$, this is equivalent to the existence of a cofinal functor $I \longrightarrow C/F$, or again $I \longrightarrow C_0/F$.)} (\footnote[66]{\textbf{N.B.} We shall denote by $\mathrm{Filt}^{\pi} \subset \Mc$ the strictly full subcategory of $\Mc$ consisting of the $F$ satisfying this condition.}).
\item[\textit{b$'$)}] \textit{We have $F \in$} $\mathrm{Ob}\,\mathrm{Filt}'^{\pi}$\textit{, where} $\mathrm{Filt}'^{\pi}$ \textit{is the smallest strictly full subcategory of $\Mc$ containing $C$, and stable under $\varinjlim$ large relative to $\pi_0$ and of cardinality $\leqslant \pi$.}
\item[\textit{c)}] \textit{We have $F \in$} $\mathrm{Ob}\,\Mc_{\pi}$\textit{, i.e. $F$ is a $\pi$-accessible object of $\Mc$.}
\item[\textit{c$'$)}] \textit{We have $F \in$} $\mathrm{Ob}\, C^{\wedge}_{\pi}$\textit{, i.e. $F$ is a $\pi$-accessible object of $C^{\wedge}$.}
\end{itemize}
\textit{This being set, we have the following diagram of implications.}
\[
\begin{tikzcd}[column sep=large, row sep=large]
\mathrm{a} \arrow[r, Rightarrow] \arrow[dr, Rightarrow] & \mathrm{a}' \arrow[r, Rightarrow] & \mathrm{a}'' \arrow[dr, Rightarrow] & & \\
& \mathrm{b} \arrow[r, Rightarrow] & \mathrm{b}' \arrow[r, Rightarrow] & \mathrm{c}' \arrow[dr, Rightarrow] \arrow[ulll, dotted, bend right=50, "\text{if } \pi \geqslant \pi_1"'] & \\
& & & & \mathrm{c}\;, \arrow[ull, dotted, bend left=30, "\text{if } \pi = \pi^{\pi_0}"']
\end{tikzcd}
\]
\textit{where the two dotted implications, c$'$ $\Longrightarrow$ a and c $\Longrightarrow$ b, are conditional subject to the condition $\pi \geqslant \pi_1$, resp. $\pi = \pi^{\pi_0}$.}

\msp{page 163}

\textbf{Corollary 4.7.2.}
\begin{itemize}[leftmargin=4em]
\item[\textit{1°)}] \textit{If $\pi \geqslant \pi_1$, the six conditions a, a$'$, a$''$, b, b$'$, c$'$ are equivalent, and the condition in question (which can be written $F \in$} $\mathrm{Ob}\,\mathrm{Filt}^{\pi}$ \textit{where $F \in$} $\mathrm{Ob}\,\mathrm{Filt}'^{\pi}$ \textit{indifferently) implies c :} % [sic] où -> ou, alternative "ou ... indifféremment"
\[
\mathrm{Filt}^{\pi} \;=\; \mathrm{Filt}'^{\pi} \;\subset\; \Mc_{\pi} \qquad \textit{if } \pi \geqslant \pi_1 \tag*{(4.7.2.1)}
\]
\textit{(whereas in general, for $\pi \geqslant \pi_0$, we have}
\[
\mathrm{Filt}^{\pi} \;\subset\; \mathrm{Filt}'^{\pi}\ (\,\subset C^{\wedge}_{\pi} \cap \Mc\,) \;\subset\; \Mc_{\pi} \qquad \textit{if } \pi \geqslant \pi_0\textit{).} \tag*{(4.7.2.2)}
\]
\item[\textit{2°)}] \textit{When $\pi = \pi^{\pi_0}$, e.g. when $\pi$ is of the form $2^{c}$, with $c \geqslant \pi_0$, then conditions b, b$'$, c, c$'$ are equivalent, i.e. we have}
\[
\mathrm{Filt}^{\pi} \;=\; \mathrm{Filt}'^{\pi} \;=\; (C^{\wedge})_{\pi} \cap \Mc \;=\; \Mc_{\pi}\;. \tag*{(4.7.3.2)} % [sic] (4.7.3.2) -> (4.7.2.3), the neighbouring labels are 4.7.2.1, 4.7.2.2, 4.7.2.4
\]
\item[\textit{3°)}] \textit{When $\pi \geqslant \pi_1$} and \textit{$\pi = \pi^{\pi_0}$, then the seven conditions a), a$'$), a$''$), b) , b$'$), c) , c$'$) are equivalent, so in particular}
\[
\mathrm{Filt}^{\pi} \;=\; \mathrm{Filt}'^{\pi} \;=\; C^{\wedge}_{\pi} \cap \Mc \;=\; \Mc_{\pi} \qquad (\pi \geqslant \pi_1 \textit{ and } \pi = \pi^{\pi_0}\textit{).} \tag*{(4.7.2.4)}
\]
\end{itemize}

\end{document}
